\section{Limitations and Future Research Directions}
\label{sec:limitations}

\subsection{Limitations of scope}

\paragraph{Drug coverage.} The residual analysis is built on a cache of $620{,}564$ cells covering
$106$ drugs across $50$ cell lines, sampled from Tahoe's shards --- roughly a tenth of the atlas's
compound set. This is ample to test whether the model uses the drug, and to measure transfer, but it
is not a population estimate. Any claim about the \emph{distribution} of drug behaviours should be
read against the broader characterisation, which used $6{,}628$ conditions.

\paragraph{Only reproducible conditions.} The residual arm retains the $62\%$ of conditions whose
residual reproduces across a half-split ($\cos > 0.2$). This is the right choice --- training on
irreproducible residuals is training on noise --- but it means results describe the winnable
fraction of the task, not all of it. The filter might be expected to select \emph{toward} context-independent conditions and so bias the transfer coefficient upward. Measured, it does the opposite: the unfiltered estimate is $0.598$ against the filtered $0.557$. The filtered value is the one quoted, because disattenuation is unstable at the very low reliabilities the filter removes.

\paragraph{The intervals omit generation variance.} Every confidence interval in this thesis is a
clustered bootstrap that resamples \emph{cell lines} while treating each condition's score as fixed.
It therefore captures between-cell-line variance and not the variance introduced by sampling the
model's output. With four generations per arm at temperature $0.8$ the second source is not
negligible: two evaluations of the same checkpoint on the same conditions, differing only in that
one carried two additional scramble arms and so consumed the sampling stream differently, returned
$\gap$ on the held-out-drug split as \ci{+0.0195}{-0.045}{+0.086} and \ci{+0.1114}{+0.051}{+0.173}
respectively. The headline arms were stable across the pair, and the held-out-drug arm --- the
smallest, at $n=120$ --- was not. \todo{seed replication in flight; report the across-seed spread
beside the within-seed interval and state which sources each captures} Until that lands the
held-out-drug arm should be read as underdetermined rather than as a null, which weakens its use as
the control licensing the transfer result.

\paragraph{Pseudobulk resolution.} Residuals are condition-level pseudobulks over $\geq 40$ treated
cells. The single-cell heterogeneity of the response is not modelled in this frame, and the
project's own earlier results show why: at single-cell resolution the drug moves a cell less than
noise does (SNR $\approx 0.75$, median $0$ differentially expressed genes at FDR $q<0.05$).

\paragraph{Cross-plate absolutes.} The residual-space scores are computed with comparison sets
grouped by cell line rather than by plate. The \emph{difference} $\gap$ is leak-immune, because both
arms share the same control cell, but the absolute $\NIR$ values in that frame are not directly
comparable to the within-plate tables elsewhere. \todo{rescore the generation-free arms within plate
--- this costs no GPU time}

\paragraph{Frame coverage.} The drug-specific residual is $\lVert\resid\rVert/\lVert\shift\rVert =
0.786$ of the response, against $0.620$ for the generic program and batch, and the two are
essentially orthogonal ($0.786^2 + 0.620^2 = 1.002$). So this frame covers $62\%$ of the response
\emph{variance} and every result here should be read within that bound --- a real but not total
share. The remaining $38\%$ is precisely the component the field's absolute-accuracy metrics are
dominated by, which is why a drug-agnostic predictor scores so well on them.

\subsection{Limitations of the unseen-drug evaluation}

Two defects, both of which must be stated rather than papered over.

\paragraph{It is not a clean held-out-drug design.} The prompt contains
\texttt{Mechanism: \{moa\}}, so holding out a drug still hands the model the held-out drug's class
label. The correct design is leave-one-\emph{mechanism}-out. The leak is probably small ---
\texttt{moa\_lookup} scores $0.529$, so the mechanism label carries little about the response --- but
the design is wrong and no unseen-drug claim should be made until it is fixed.

\paragraph{It is underpowered by construction.} The confidence interval half-width on the
\texttt{unseen\_drug} gap is $\approx \pm 0.087$, against a maximum effect available through the
mechanism channel of $\approx 0.033$. The arm therefore \emph{cannot} detect the largest effect it
could plausibly be measuring. Its null is a statement about power, not about absence, and is
reported as such.

\subsection{Directions}

\paragraph{The arm this work closed rather than opened.} The obvious continuation of
Section~\ref{sec:res-scope} is a channel-conditioned prompt: append protein targets to the
instruction, retrain, and evaluate on a leave-one-mechanism-out split. We do not recommend it, and
the reason is measured rather than argued. Section~\ref{sec:res-field} shows the model reads the drug
\emph{name} (\ci{+0.081}{+0.059}{+0.101}) and does not measurably read the \emph{mechanism} field it
is already given (\ci{+0.009}{-0.017}{+0.038}), in an arm powered to detect an effect the size of the
name's. Adding a target field to a prompt whose mechanism field is ignored addresses information
availability, and the binding constraint is the readout.

What would be worth building instead is an intervention that makes the drug \emph{causal} rather than
merely present --- a learned per-drug modulation of the transformation from control to treated, so
that the drug parameterises the computation instead of sitting in the residual stream as an input the
readout can round away. That is a research architecture rather than a prompt change, and it is
outside the scope of what remained here.

\paragraph{Two channel questions that do remain open.} Target and mechanism overlap, so whether
target adds anything \emph{over} mechanism is unresolved --- and it matters, because the prompt
already exposes mechanism and the model already ignores it. Separately, the channel gate measures
\emph{similarity transfer}: an unseen drug resembling a seen drug with the same target. It does not
test a \emph{compositional rule} --- that a named target's own pathway moves, which would apply to
any drug whose target is named, including one with no similar drug in training. The second is the
stronger claim, it does not reduce to a lookup, and it is untested.

\paragraph{Unseen cell lines.} A cell line never seen with any drug is a different generalisation
axis from the held-out pairings tested here, and one not blocked by the structure gate: it asks
whether the control cell's state is used \emph{compositionally} with the drug. The transfer
coefficient now scopes it: with the interaction both substantial ($\approx45\%$) and idiosyncratic,
a new cell line offers no consistent direction to exploit, so the achievable target there reduces to
recalling the drug's per-drug constant --- which a lookup already does at $0.963$. The axis is
therefore less promising than it appeared before Section~\ref{sec:res-transfer} was measured.

\paragraph{Objectives that were considered and rejected.} Recorded so they are not re-litigated:
\begin{itemize}
  \item \textbf{Reinforcement learning with a contrastive reward.} Its mechanism --- a per-sequence
        reward, so the token-dilution ratio never enters the gradient --- attacks a bottleneck that
        the residual encoding has already lifted ($34.6 \rightarrow 118.2$ differing tokens); and its
        reward optimises similarity to the true residual, which is precisely what
        \texttt{drug\_lookup} already achieves at $0.963$. Its best case is convergence to the
        lookup.
  \item \textbf{A differentiable profile-level contrastive loss.} The natural construction,
        $\textrm{score}(g) = \sum_i P(\textrm{token}_i = g)\,w(i)$, presumes gene $\approx$ token,
        but $86.9\%$ of the $946$ panel genes share a first sub-token under this tokeniser (mean
        $3.254$ BPE tokens per gene, only 8 single-token). Under teacher forcing the difficulty
        compounds: the residual target \emph{is} the drug signature, so the true prefix identifies
        the drug within a handful of symbols and the gradient vanishes exactly at the
        drug-blind solution the loss exists to penalise.
  \item \textbf{Loss re-weighting toward differentially expressed genes.} Re-weighted cross-entropy
        against a single target contains no negatives, so nothing in it forces the output for one
        drug to differ from another's. Given that the readout is direction-blind
        (Section~\ref{sec:res-mechanism}), it addresses token dilution but not the actual defect.
\end{itemize}

\paragraph{Generality.} Everything here is one backbone on one atlas in one representation. Whether
the readout-specificity failure is a property of cell-sentence models specifically, or of
autoregressive conditional generation over transcriptomes generally, is open. The mechanistic
protocol of Section~\ref{sec:methods-probe} --- purified subspace, removal gate, positive control,
matched-norm swap --- transfers unchanged to any model with an accessible residual stream, and is
probably the most reusable artifact of this work.
